\section{复现材料与数据生成}
\label{app:repro}

论文源码位于 \path{paper/lasa_journal_cn/}。执行
\code{build.ps1} 会先运行证据生成器，再编译 PDF。生成器直接读取：
\begin{itemize}[leftmargin=1.8em]
  \item \path{build_abi_v2_release_r5/reports/system_impl_gate.txt}；
  \item \path{tools/coco80/model_spec.json}；
  \item 正式性能 \path{performance_final_summary/summary.json}；
  \item 部署验证 \path{verification_snapshot.json}；
  \item 外部 COCO 评测目录不可用时，显式使用
  \path{data/evidence_snapshot.json}。
\end{itemize}
每个输入和生成的 tex 表均记录 SHA256 于
\path{generated/evidence_manifest.json}。表中 PENDING 是数据状态而不是
LaTeX 占位符；冻结脚本必须拒绝仍含 PENDING 的投稿 PDF。

公开 artifact 计划包含 RTL、Vitis 软件、测试/构建脚本、论文表生成器、
r5 BIT/XSA/ELF、参数 manifest、固定小型 golden 和结果摘要。COCO 原图和
annotations 不重新分发，只提供官方获取说明、image-id 列表与 SHA256。
大型 raw node/head 文件可以放在外部对象存储，仓库保留 index、大小、hash
和下载验证工具。

复现流程分为“数据重算”和“投稿冻结”两个层级。数据重算允许外部大文件暂时
不可用，此时必须显式读取带来源说明的 snapshot，并在页首显示证据类型；它
只用于继续编辑内部稿。投稿冻结则要求所有原始报告可访问、生成表格的哈希
与 manifest 一致、九张核心图和参考文献数量满足合同，并且六项 publication
gate 全部为 PASS。严格检查发现任一 PENDING、来源漂移、手工正文表格或摘要
禁用表述都会返回非零状态，不生成“可投稿”标志。这样，即使 PDF 可以正常
编译，也不会被误认为实验已经完整。

最终 release tag 还应保存生成环境版本、Git 状态和构建命令，并在一台没有
本地缓存的新环境执行自检。COCO 原图不进入仓库，但 image-id 顺序、下载来源、
annotation 哈希和预处理合同必须足以重建相同输入；受版权或容量限制放在外部
的 raw tensor 也必须先校验大小与 SHA256 再参与比较。论文中的每个数值由
JSON/CSV 和生成脚本产生，人工只撰写解释文字，因此模型参数替换或新消融
完成后，可以重建表格并自动暴露哪些结论需要同步修改。

\begin{table*}[t]
\centering
\caption{主要 artifact 身份。完整 64 位哈希保存在构建/结果 manifest 中。}
\label{tab:artifact-id}
\small
\begin{tabular}{lll}
\toprule
Artifact & SHA256 前后缀 & 作用\\
\midrule
r5 BIT & 见证据 manifest & XCK26 200 MHz 配置 \\
r5 XSA & 见证据 manifest & Vitis 平台与硬件描述 \\
上游 YOLOv3-tiny 权重 & 74fb61c9...946280 & COCO80 FP32 资产 \\
epoch-1 参数包 & 974970bb...3f11b & 当前功能部署 INT8 \\
正式性能 summary & 40bdccc5...d7aa & 3x1000 统计 \\
600 s soak summary & fa22bb62...987d & 稳定性证据 \\
\bottomrule
\end{tabular}
\end{table*}

